\section{Hilo conductor de la clase: Frontier Product Research no es «terminar la investigación y luego empaquetarla»}

El título de esta clase puede malinterpretarse fácilmente como una introducción común al aprendizaje por refuerzo (Reinforcement Learning, RL). Lo que Karina Nguyen discute en realidad es un problema de sistemas más concreto: cuando el modelo entra en la escritura, el software, la colaboración, la educación y las tareas de largo plazo, \textbf{cómo quiere el producto que actúe el modelo}, \textbf{cómo convierte la investigación esa conducta en entornos y recompensas} y \textbf{cómo permite la interfaz que el usuario saque a la luz los fallos} deben diseñarse en conjunto dentro de un mismo ciclo cerrado. Si la investigación optimiza primero un benchmark desvinculado del producto y después entrega el modelo al equipo de producto para empaquetarlo, muchos fallos reales no aparecen sino hasta el lanzamiento.

\lecturefigure{slide-01-golden-age-curiosity.jpg}{La edad de oro de «la curiosidad, la creación de conocimiento y la fabricación»: la conferencia plantea primero una orientación de valores}{00:00:49--00:01:41}

\teachervoice{00:00:49--00:01:38, Karina aclara que esto no es una lecture unidireccional, sino que espera una discusión colaborativa en el aula; su punto de partida es que cada persona puede construir algo significativo con ayuda de la IA. Este juicio de valor explica por qué lo que sigue discute a la vez creativity, agency, verification y alignment.}

\begin{importantbox}{Definición: Frontier Product Research}
En esta clase, \term{frontier product research} no significa «seguir los modelos más recientes y hacer productos». Toma las capacidades todavía inestables del modelo como material de investigación, convierte las tareas reales de los usuarios en entornos de entrenamiento y evaluaciones, y luego usa la interacción con el producto para reunir contraejemplos con rapidez. Su ciclo cerrado mínimo es
\[
\begin{aligned}
\text{product belief}&\rightarrow \text{task/env}\rightarrow \text{reward/eval}\\
&\rightarrow \text{model behavior}\rightarrow \text{user evidence}\rightarrow \text{next hypothesis}.
\end{aligned}
\]
Aquí, product belief es el principio de conducta que se espera que el producto muestre al final; debe traducirse en tareas observables y no quedarse en un eslogan.
\end{importantbox}

\begin{knowledgebox}{Límites de la evidencia en esta clase}
La primera mitad usa muchas demos de producto. Estas imágenes pueden demostrar que cierta interacción funcionaba en ese momento y mostrar cómo entiende el ponente la dirección del producto, pero no bastan por sí solas para demostrar que el modelo es confiable para todos los usuarios, distribuciones y niveles de riesgo. Lo que sigue (eval, la taxonomy de rechazos, XSTest, RL environment y reward hacking) cubre precisamente las preguntas que una demo no puede responder.
\end{knowledgebox}

\subsection{Canvas: un cambio de interfaz no es decoración, sino un cambio en la estructura de la tarea}

El primer ejemplo descompone una solicitud educativa en explicación, generación de código, ejecución y visualización. Si el sistema solo ofrece un mensaje de chat, al usuario le cuesta leer el concepto, modificar el código y observar la gráfica al mismo tiempo; Canvas eleva la «conversación» a un espacio de trabajo compartido, de modo que la salida del modelo se vuelve un objeto editable. Al leer la figura, no basta con mirar la curva Gaussian: hay que ver cómo la explicación de la izquierda, el código del centro y el resultado de la ejecución de la derecha forman un ciclo verificable.

\lecturefigure{slide-02-gaussian-canvas.jpg}{La explicación de la distribución Gaussian, el código en Python y el resultado de su ejecución se colocan en el mismo espacio de trabajo}{00:02:05--00:02:57}

Para una variable aleatoria $X\sim\mathcal{N}(\mu,\sigma^2)$, la densidad es
\[
p(x)=\frac{1}{\sqrt{2\pi}\sigma}
\exp\!\left(-\frac{(x-\mu)^2}{2\sigma^2}\right).
\]
Aquí $\mu$ controla el centro y $\sigma$ controla la escala. Que el modelo dé la fórmula no significa que la enseñanza haya terminado; los símbolos abstractos se convierten en una comprensión operable solo cuando el usuario puede modificar los parámetros, ejecutar el código y observar cómo cambia la curva. Por eso el diseño del producto cambia la pregunta de evaluación: ya no se pregunta solo si la respuesta es correcta, sino también si el código se ejecuta, si la gráfica corresponde a la explicación y si las modificaciones surten efecto de forma local.

\teachervoice{00:02:05--00:04:04, el ponente usa la enseñanza de la Gaussian y capturas de pantalla de artículos para explicar la motivación de Canvas: ChatGPT era al principio una conversational UI, pero el código, los textos largos y las preguntas de seguimiento detalladas expusieron las limitaciones de la interfaz de chat. Nota de clase: el cambio de UI proviene de la estructura de la tarea, no de un rediseño visual.}

El ejemplo de comprensión de artículos muestra además un «punto de entrada para la verificación». El usuario proporciona primero la gráfica de calibration del artículo y luego pide al modelo que la explique y que genere un informe de investigación y código. El material original, la explicación del modelo y los nuevos productos deben poder rastrearse lado a lado; de lo contrario, el sistema tiende a confundir una narración fluida con un hecho.

\lecturefigure{slide-03-paper-calibration-source.jpg}{La gráfica de calibration del artículo: la explicación del modelo debe remitir a la evidencia visual original}{00:02:57--00:03:09}

\lecturefigure{slide-04-calibration-canvas-report.jpg}{El informe de calibration generado en Canvas: se selecciona una parte del contenido y se continúa preguntando}{00:03:09--00:04:12}

\begin{knowledgebox}{Primer uso: Calibration}
\term{Calibration (calibración)} describe si la confianza del modelo coincide con su tasa real de aciertos. Si, entre todas las predicciones con confianza cercana a $q$, la proporción real de aciertos también es cercana a $q$, ese intervalo está aproximadamente calibrado. Un error frecuente es confundir calibration con accuracy: un modelo puede ser preciso pero exceder su confianza, o bien ser poco preciso pero estar calibrado en sentido probabilístico. El producto necesita mostrar por separado «cuál es la respuesta» y «cuánto debe creerse».
\end{knowledgebox}

\lecturefigure{slide-05-claude-usage-bachelor-degrees.jpg}{Gráfica de correlación de un informe de uso educativo: los datos del producto pueden generar preguntas de investigación}{00:04:12--00:04:45}

Esta gráfica de uso relaciona los usos educativos de Claude.ai con distintas carreras de licenciatura en Estados Unidos. Para leer la figura, primero se observa en el eje horizontal la proporción de cada carrera o la proporción de uso, y luego la posición de los puntos que se apartan de la línea de tendencia; la gráfica puede sugerir que las distintas disciplinas tienen patrones de adopción diferentes, pero no puede demostrar que «la formación en una carrera causa el uso» ni que «el uso causa resultados de aprendizaje». La telemetry del producto puede generar hipótesis que vale la pena investigar, pero la correlación, el sesgo de selección y el autorreporte de los usuarios deben tratarse por separado.

\begin{warningbox}{Tres filtros entre los registros del producto y una conclusión científica}
Primero, los usuarios no son una muestra aleatoria de la población; segundo, el número de clics o de generaciones no equivale al éxito en la tarea; tercero, una correlación a nivel de grupo no puede explicar directamente la causalidad individual. Frontier product research puede detectar problemas rápidamente a partir del tráfico real, pero sigue siendo necesario precisar el sampling, el outcome y el counterfactual.
\end{warningbox}

\subsection{Resumen de la sección}

La apertura de la clase ya fija la escala de toda la sesión: las capacidades del modelo, la interfaz de trabajo, la evidencia de los usuarios y los objetivos del posentrenamiento forman un solo sistema. El valor de Canvas no es «tener una ventana más», sino reunir la explicación, la edición, la ejecución y la verificación en una misma interacción; el valor de los datos del producto no es que sean verdaderos por naturaleza, sino que vuelven más concretas las preguntas de investigación.
